\section{Related Work}
\label{sec:related_work}

\subsection{LLM for EDA Scripting}
\label{sec:rw_eda}

LLMs for chip design have largely targeted RTL generation: ChipNeMo, RTLCoder, AutoChip, and CodeV-R1~\cite{liu2023chipnemo,liu2024rtlcoder,thakur2023autochip,zhu2026qimeng} domain-adapt or fine-tune models for Verilog, relying on human-constructed testbenches, which is orthogonal to operating an EDA tool's API surface. The procedural-scripting strand has grown quickly: ChatEDA, ChaTCL, EDA-Copilot, iScript, FluxEDA, JARVIS, and AutoEDA~\cite{wu2024chateda,rui2025chatcl,xiao2025eda,xu2026iscript,chen2026fluxeda,lu2025autoeda,zhang2025jarvis} span flow planning, Tcl/Q\&A assistance, multi-agent collaboration, single-turn generation, and execution infrastructure. Within this strand, ZhuLong~\cite{liu2026zhulong} is the direct precursor to our work: it combines API retrieval, documentation inspection, sandbox execution, and offline API self-exploration, but solves tasks independently without reusing execution traces across tasks. Another closely related system is OpenROAD Agent~\cite{wu2025openroad}, which adds execution-feedback self-correction on an open physical-design tool, but fine-tunes the model and does not externalize cross-task experience. Our scenario shares its generate-execute-revise shape yet keeps the model frozen and judges pass/fail by real sandbox execution with hidden tests---addressing cross-task reuse through self-evolution rather than fine-tuning. The Agentic EDA survey~\cite{zang2025dawn} situates this line at L3 autonomy. Two gaps recur for a frozen-agent, hidden-test setting: none externalizes cross-task experience into a reliable, leakage-controlled store, and none carries a domain-level static check beyond the tool's own schema---both of which our system addresses.
% 面向芯片设计的LLM主要面向RTL生成：ChipNeMo、RTLCoder、AutoChip和CodeV-R1~\cite{liu2023chipnemo,liu2024rtlcoder,thakur2023autochip,zhu2026qimeng}通过领域适配或微调模型生成Verilog，依赖人工构建的测试台，这与操作EDA工具API层面正交。流程脚本编写方向发展迅速：ChatEDA、ChaTCL、EDA-Copilot、iScript、FluxEDA、JARVIS和AutoEDA~\cite{wu2024chateda,rui2025chatcl,xiao2025eda,xu2026iscript,chen2026fluxeda,lu2025autoeda,zhang2025jarvis}涵盖流程规划、Tcl/问答辅助、多智能体协作、单轮生成和执行基础设施。在这一方向中，ZhuLong~\cite{liu2026zhulong}是本文最直接的前置工作：它结合API检索、文档查阅、沙箱执行和离线API自探索，但各任务独立求解，不复用跨任务执行轨迹。另一项密切相关的系统是OpenROAD Agent~\cite{wu2025openroad}，它在开源物理设计工具上增加执行反馈自修正能力，但会微调模型且未外部化跨任务经验。我们的场景与其共享生成-执行-修正形态，但保持模型冻结，并通过真实沙箱执行和隐藏测试判断通过/失败——通过自进化而非微调实现跨任务重用。Agentic EDA综述~\cite{zang2025dawn}将这一方向定位为L3自主级别。对于冻结智能体、隐藏测试的场景，仍有两个缺口：现有工作没有将跨任务经验外部化到可靠且受泄漏控制的存储中，也没有提供超越工具自身模式的领域级静态检查——本文系统同时解决这两个问题。

\subsection{Self-Evolving Agents with Governed Memory}
\label{sec:rw_self_evolving}

A broad family of works equips agents with reusable external memory: skill libraries, cross-task insight distillation, hierarchical memory, and cognitive architectures~\cite{wang2023voyager,zhao2024expel,packer2023memgpt,sumers2023cognitive}, with CodeAct~\cite{wang2024executable} unifying the action space as executable Python. The survey by Zhang et al.~\cite{zhang2024memory} provides a comprehensive taxonomy of memory mechanisms for LLM agents. These works, however, generally assume an open world or a train/test-disjoint split, which sidesteps the same-task regression signal that our fixed-set setting must surface. DemoEvolve~\cite{demoevolve2026} explicitly studies sparse terminal feedback in harness evolution, but addresses it through human demonstrations—whereas our setting assumes no such external guidance.
% 大量工作为智能体配备了可重用的外部记忆：技能库、跨任务洞察蒸馏、分层记忆和认知架构~\cite{wang2023voyager,zhao2024expel,packer2023memgpt,sumers2023cognitive}，其中CodeAct~\cite{wang2024executable}将动作空间统一为可执行的Python。Zhang等人的综述~\cite{zhang2024memory}为LLM智能体的记忆机制提供了全面的分类体系。然而，这些工作通常假设开放世界或训练/测试不相交的划分，这规避了我们的固定集合设定必须面对的同一任务回归信号问题。DemoEvolve~\cite{demoevolve2026}明确研究了框架进化中的稀疏终端反馈，但通过人类示范来解决——而我们的设定假设不存在此类外部指导。

Treating the skill store as an optimization target, SkillOpt~\cite{yang2026skillopt} edits a single skill document under a held-out gate and observes that ``plausible textual diagnoses can still hurt,'' echoing our concern that unverified experience can mislead. SkillOS~\cite{ouyang2026skillos} learns curation via RL with a four-dimensional judge, and Trace2Skill has been explored in two complementary forms: Ni et al.~\cite{ni2026trace2skill} propose a general framework that distills execution trajectories into reusable skills through parallel trajectory analysis and inductive pattern consolidation, using a verifier primarily for quality filtering of the extracted skills; Du and Pinckney~\cite{du2026trace2skill} present a test-time scaling framework for hardware agents where dense verifier feedback actively drives the oracle-mutator-selector evolution loop. These works optimize on disjoint splits with relatively dense signal, whereas our sparse-pass/fail, fixed-set setting makes which-experience-to-trust more pressing than how to edit a document.
% 将技能存储视为优化目标，SkillOpt~\cite{yang2026skillopt}在留出门控下编辑单个技能文档，并观察到“看似合理的文本诊断仍然可能造成损害”，这与我们对未验证经验可能误导的担忧一致。SkillOS~\cite{ouyang2026skillos}通过四维评判器的强化学习来学习策展，Trace2Skill已从两个互补方向进行探索：Ni等人~\cite{ni2026trace2skill}提出了一个通用框架，通过并行轨迹分析和归纳模式整合将执行轨迹蒸馏为可复用技能，其中验证器主要用于蒸馏后技能的质量过滤；Du和Pinckney~\cite{du2026trace2skill}提出了一个面向硬件智能体的测试时扩展框架，其中密集验证器反馈主动驱动"预言机-变异器-选择器"进化循环。这些工作在不相交划分上优化，且信号相对密集，而我们的稀疏通过/失败、固定集合设定使得“信任哪些经验”比“如何编辑文档”更为紧迫。

Verifier-guided evolution brings the reward signal closer to our hidden-test setting. ReVeal~\cite{jin2025reveal} self-generates test cases for dense rewards, but self-generated tests can be gamed and do not substitute for a hidden verifier. Du and Pinckney~\cite{du2026trace2skill} share our hidden-verifier, no-answer framing but run online per task with physical isolation and task-wise skills; their finding that cross-task transfer is unreliable for heterogeneous RTL cautions how far logical-error knowledge should travel—a concern we address through governed memory and anti-leakage retrieval. EvoAgentBench~\cite{gao2026evoagentbench}, the most comparable benchmark, formalizes this tension with an ability-supported yet instance-disjoint split, finding that an Anchor Skill oracle helps while no automatic method sustains the gain.
% 验证器引导的进化使奖励信号更接近我们的隐藏测试设定。ReVeal~\cite{jin2025reveal}自我生成测试用例以获得密集奖励，但自我生成的测试可能被操纵，不能替代隐藏验证器。Du和Pinckney~\cite{du2026trace2skill}与我们共享隐藏验证器、无答案的框架，但在每个任务上在线运行，具有物理隔离和任务级技能；他们发现跨任务迁移在异构RTL上不可靠，这警示了逻辑错误知识应传播多远——我们通过受治理记忆和反泄漏检索来应对这一问题。EvoAgentBench~\cite{gao2026evoagentbench}是最具可比性的基准测试，通过能力支撑但实例不相交的划分形式化了这种张力，发现Anchor Skill预言机有帮助，但没有任何自动方法能维持其收益。

Several works inform the governance machinery underlying our approach: MUSE-Autoskill~\cite{lin2026muse} on refinement and reflection, Externalization~\cite{zhou2026externalization} and Reflexion~\cite{shinn2023reflexion} on memory artifacts and same-task iteration, and Governing Evolving Memory~\cite{lam2026governing} on System-2 verification. LastingBench~\cite{fang2025lastingbench} proposes counterfactual rewriting to safeguard benchmarks against knowledge leakage, underscoring the importance of our anti-leakage retrieval design. The central lesson is that evolving memory requires an external gate grounded in ground truth to remain reliable. Our method addresses this challenge under sparse pass/fail supervision and strict anti-leakage constraints.
% 若干工作为我们方法背后的治理机制提供了参考：MUSE-Autoskill~\cite{lin2026muse}涉及精炼和反思，Externalization~\cite{zhou2026externalization}和Reflexion~\cite{shinn2023reflexion}涉及记忆制品和同任务迭代，Governing Evolving Memory~\cite{lam2026governing}涉及系统2验证。LastingBench~\cite{fang2025lastingbench}提出通过反事实重写来保护基准免受知识泄露，凸显了我们反泄漏检索设计的重要性。反复出现的教训是：进化的记忆需要一个外部的、真值锚定的门控来保持可信赖——这在稀疏通过/失败和反泄漏约束下，正是我们的方法所解决的问题。
